\documentclass[8pt, a4paper]{article}

% ── PACKAGES ──────────────────────────────────────────────────
\usepackage[utf8]{inputenc}
\usepackage[T1]{fontenc}
\usepackage{geometry}
\geometry{top=0.6cm, bottom=0.6cm, left=1.0cm, right=1.0cm}
\usepackage{hyperref}
\hypersetup{
  colorlinks=true, urlcolor=blue, linkcolor=blue,
  pdftitle={Plant Architecture Inversion --- One Page Summary v2.3},
  pdfauthor={Eric Robert Lawson / OrganismCore}
}
\usepackage{microtype}
\usepackage{parskip}
\usepackage{booktabs}
\usepackage{array}
\usepackage{xcolor}
\usepackage{mdframed}
\usepackage{multicol}
\usepackage{titlesec}
\usepackage{lmodern}

% ── COLOURS ───────────────────────────────────────────────────
\definecolor{headergreen}{RGB}{20,100,60}
\definecolor{statementbox}{RGB}{20,100,60}
\definecolor{rulecolor}{RGB}{20,100,60}

% ── SECTION FORMAT ────────────────────────────────────────────
\titleformat{\section}
  {\small\bfseries\color{headergreen}}
  {}{0em}{}[\vspace{-0.3em}
  \textcolor{rulecolor}{\rule{\linewidth}{0.4pt}}]
\titlespacing{\section}{0pt}{0.2em}{0.05em}

\setlength{\parskip}{0.05em}
\setlength{\parindent}{0pt}
\setlength{\columnsep}{0.7em}
\pagestyle{empty}

% ══════════════════════════════════════════════════════════════
\begin{document}

% ── HEADER ────────────────────────────────────────────────────
\begin{center}
  {\large\bfseries\color{headergreen}
  THE PLANT ARCHITECTURE INVERSION EXPERIMENT}\\[0.03em]
  {\small\color{headergreen}
  Coherence Gradient Field Specification Produces
  Pre-Specified Architectural Inversion
  in \textit{Raphanus sativus} --- Operative Protocol v2.3}\\[0.03em]
  {\footnotesize
  Eric Robert Lawson\ $\cdot$\
  OrganismCore\ $\cdot$\
  ORCID: \href{https://orcid.org/0009-0002-0414-6544}{0009-0002-0414-6544}\ $\cdot$\
  \href{mailto:OrganismCore@proton.me}{OrganismCore@proton.me}\ $\cdot$\
  2026-03-18}
\end{center}

\vspace{0.03em}

% ── CLAIM BOX ─────────────────────────���───────────────────────
\begin{mdframed}[
  backgroundcolor=statementbox, linecolor=statementbox,
  linewidth=0pt,
  innertopmargin=0.2em, innerbottommargin=0.2em,
  innerleftmargin=0.6em, innerrightmargin=0.6em
]
\begin{center}
\footnotesize\bfseries\color{white}
The organism is not the output of a genome alone.
It is the output of a genome operating within a coherence gradient field.
Change the field --- the genome executes a different program.
Same genome. Different field. Different organism.\\
Placing a radish seed in an \textbf{inverted coherence gradient field}
--- hydroponic foam above, UV-A embedded foam below,
complete external light deprivation ---
produces \textbf{root growing upward, shoot growing downward.}
Derived geometrically. \textbf{Prediction locked before the experiment.}
Pre-registered:
\href{https://doi.org/10.5281/zenodo.19040399}{\texttt{doi.org/10.5281/zenodo.19040399}}
\end{center}
\end{mdframed}

\vspace{0.1em}

% ── TWO COLUMNS ───────────────────────────────────────────────
\begin{multicols}{2}

% ── LEFT COLUMN ───────────────────────────────────────────────

\section{The Theorem}
{\small
\textbf{Standard:} Genome $\rightarrow$ Organism.\quad
\textbf{Complete:} Genome $+$ Field $\rightarrow$ Organism.

Standard architecture (root down, shoot up) is the attractor
the genome resolves to in the standard gradient field.
\textbf{Invert the field before germination: the genome resolves
to the inverted attractor.} No genetic modification. One generation.

\smallskip\noindent
\textbf{Root} navigates toward water and minerals.
\textbf{Shoot} navigates toward the only photonic signal in the
field --- UV-A embedded within the substrate below.
}

\section{Signal Deprivation Principle}
{\small
Scenarios 2--5 are in \textbf{complete external light deprivation.}
The embedded UV-A LED is the \textit{only} photonic signal in the
coherence field. The seedling is in \textbf{skotomorphogenic state}
--- maximum photoreceptor sensitivity, zero ambient noise.

\textit{A book falling in a silent room is orders of magnitude
louder than a shout in a noisy one.}

UV-A efficiency relative to peak blue light is irrelevant.
Signal-to-noise ratio is effectively infinite.
Light exclusion is the load-bearing engineering constraint.
}

\section{Pre-Registered Prediction --- Scenario 3}
{\small
\textbf{Root:} Upward. Hydrotropism overrides gravitropism.
MIZ1-mediated amyloplast degradation suppresses gravity sensing
(Zhang et al.\ PNAS 2025). Water in S1 Type~A foam above is the
only hydric signal.

\textbf{Shoot:} Downward. UV-A embedded in S2 Type~B foam below
is the only photonic signal. Shoot penetrates foam to reach it.

\textbf{AIS\,=\,2} at T$=$72\,h, maintained T$=$168\,h.\\
{\footnotesize AIS\,=\,2: root $>$135\textdegree, shoot $<$45\textdegree
from vertical. Scenario 1 must record AIS\,=\,0 in same run.
Minimum 3 seeds, 2 independent runs for reportable confirmation.}
}

\section{Confirmed Literature Support}
{\footnotesize
\begin{tabular}{p{3.6cm}p{2.0cm}}
Root upward to water, radish & Takahashi 2003 \\
MIZ1 gravity suppression, drought & Zhang PNAS 2025 \\
Hydrotropism $>$ gravitropism & PNAS 2025, 1994 \\
Shoot bends toward buried light & Phototropin lit. \\
Root $+$ve phototropism, no gravity & Kiss ISS 2012 \\
Photosynthesis from buried geothermal light & Beatty PNAS 2005 \\
Radiotropism toward radiation source & Dadachova 2007 \\
\end{tabular}

\vspace{0.1em}
No mechanism is novel in isolation. The combination,
geometric pre-specification, and simultaneous full architectural
inversion from a locked prediction are novel.
Every prior result is post-hoc or partial.
}

\columnbreak

% ── RIGHT COLUMN ──────────────────────────────────────────────

\section{The Five Scenarios --- v2.3}
{\footnotesize
\textbf{Substrate types:}
\textbf{Free UV} --- LED free in S1 air, zero penetration cost.
\textbf{Type A} --- 20\,ppi foam + hydroponic (water + minerals, no UV).
\textbf{Type B} --- 20\,ppi foam + UV-A LED embedded (UV only, no hydroponic).
\textbf{Type C} --- 20\,ppi foam + hydroponic + UV-A LED embedded (all co-located).

\vspace{0.2em}
\begin{tabular}{p{0.5cm}p{1.85cm}p{1.85cm}p{0.85cm}}
\toprule
\textbf{ID} & \textbf{S1 top} & \textbf{S2 bottom} & \textbf{Cone} \\
\midrule
S1 & Free UV in air & Type A & Apex up \\[0.1em]
S2 & Type B (UV only) & Type A & Apex up \\[0.1em]
\textbf{S3} & \textbf{Type A} & \textbf{Type B (UV only)} & \textbf{Apex dn} \\[0.1em]
S4 & Type C (all) & Empty void & Apex up \\[0.1em]
S5 & Empty void & Type C (all) & Apex dn \\
\bottomrule
\end{tabular}

\vspace{0.1em}
\textbf{Cone rule:} apex always toward void or least active substrate.\\
S3 is the primary claim. S1 is the validity gate.
S2 vs S3 is the primary inversion pair.
S4 and S5 are single-polarity maximum field tests.
}

\section{What Confirmation Establishes}
{\small
\textbf{1. Field authorship of architecture.}
Same genome, inverted field, inverted organism.
Architecture is a field output. First demonstration from a
pre-specified locked prediction.

\textbf{2. Attractor geometry is real and navigable.}
The field is the instruction set. The genome is the build capacity.
The organism is the output of both.

\textbf{3. No genome editing required.}
Specify the field. The genome executes.
}

\section{Materials and Timeline}
{\footnotesize
\begin{tabular}{p{3.3cm}p{0.75cm}p{1.5cm}}
\toprule
\textbf{Item} & \textbf{Cost} & \textbf{Source} \\
\midrule
Radish seeds & \$2--4 & Garden store \\
Reticulated foam 20\,ppi & \$5--10 & Aquarium supply \\
Hydroponic mineral solution & \$5--12 & Garden / hydro \\
Opaque PVC pipe + end caps & \$15--25 & Hardware store \\
UV-A LED strip $\sim$365\,nm & \$8--15 & Amazon \\
Metallized Mylar + funnels + net cups & \$5--11 & Dollar / kitchen \\
Aquarium silicone sealant & \$5--8 & Aquarium store \\
Red-light torch & \$5--15 & Astronomy supply \\
\midrule
\textbf{Total} & \textbf{\$50--100} & \\
\bottomrule
\end{tabular}

\vspace{0.1em}
No BSL rating. No credentials. No specialist equipment.\\
Light exclusion is the critical build requirement.\\
Primary result T\,=\,72\,h. Full confirmation T\,=\,168\,h.
}

\end{multicols}

\vspace{0.03em}
\noindent\textcolor{rulecolor}{\rule{\linewidth}{0.4pt}}

% ── FOOTER ────────────────────────────────────────────────────
\begin{center}
\footnotesize
\textit{The field came first.\ \
The genome was always listening.\ \
The architecture was never fixed.\ \
The inversion was always possible.}\\[0.1em]
\href{https://github.com/Eric-Robert-Lawson/attractor-oncology}
{\texttt{github.com/Eric-Robert-Lawson/attractor-oncology}}
\ $\cdot$\
\href{https://doi.org/10.5281/zenodo.19040399}
{\texttt{doi.org/10.5281/zenodo.19040399}}
\ $\cdot$\
\href{https://orcid.org/0009-0002-0414-6544}
{\texttt{ORCID: 0009-0002-0414-6544}}
\end{center}

\end{document}
